% ============================================================================
% WORKSHEET - Section 1.6: Applying Proportional Reasoning
% Oklahoma modified topic
% Modified States: Oklahoma
% ============================================================================
\section{Applying Proportional Reasoning}
\setcounter{wsNumber}{0}

\begin{worksheetSkills}
\begin{itemize}
    \item Estimate before solving a proportional problem exactly
    \item Check whether the exact answer is reasonable compared with the estimate
    \item Choose and justify a proportional reasoning strategy in real-world contexts
\end{itemize}
\end{worksheetSkills}

\worksheetHeader{Estimate First, Then Solve}

\begin{speedDrill}{Give an Estimate and an Exact Answer}
\timerBadge{5}

\renewcommand{\arraystretch}{1.9}
\begin{center}
\begin{tabular}{|C{6.5cm}|C{2.5cm}|C{2.5cm}|}
\hline
\textbf{Problem} & \textbf{Estimate} & \textbf{Exact} \\
\hline
$210$ miles on $6$ gallons; how far on $10$ gallons? & \answerBlank[2cm] & \answerBlank[2cm] \\
\hline
$3$ cm equals $24$ km; how far is $8$ cm? & \answerBlank[2cm] & \answerBlank[2cm] \\
\hline
$120$ pages in $8$ min; how long for $450$ pages? & \answerBlank[2cm] & \answerBlank[2cm] \\
\hline
$5$ shirts cost $\$42.50$; how much for $8$ shirts? & \answerBlank[2cm] & \answerBlank[2cm] \\
\hline
\end{tabular}
\end{center}
\end{speedDrill}
\worksheetAnswer{(1)~Estimate: about $350$ miles because $210 \div 6 = 35$ and $35 \times 10 = 350$. Exact: $350$ miles. (2)~Estimate: about $64$ km because $24 \div 3 = 8$ and $8 \times 8 = 64$. Exact: $64$ km. (3)~Estimate: about $30$ minutes because $120 \div 8 = 15$ pages per minute and $450 \div 15 = 30$. Exact: $30$ minutes. (4)~Estimate: about $\$64$ to $\$72$ because each shirt is a little over $\$8$. Exact: $42.50 \div 5 = 8.50$ dollars per shirt and $8.50 \times 8 = \$68$. Each exact answer is close to or equal to the estimate, so the results are reasonable.}

\worksheetHeader{Estimate or Exact?}

\begin{matchIt}{Match the Problem to a Smart Estimate}

\matchRow{$5$ shirts cost $\$42.50$; what about $8$ shirts?}{About $\$64$ to $\$72$}
\matchRow{$3$ cm represents $24$ km; what about $8$ cm?}{About $64$ km}
\matchRow{$\frac{3}{2}$ cups of rice feed $4$ people; what about $10$ people?}{About $3.5$ to $4$ cups}
\matchRow{$120$ pages in $8$ min; what about $450$ pages?}{About $30$ minutes}
\matchRow{Small paint is about $\$24$ per gallon; large paint is about $\$22$ per gallon}{Large can is a little cheaper}
\end{matchIt}
\worksheetAnswer{$5$ shirts cost $\$42.50$ matches ``About $\$64$ to $\$72$'' because each shirt is a little more than $\$8$. The map problem matches ``About $64$ km'' because each centimeter is about $8$ km. The rice problem matches ``About $3.5$ to $4$ cups'' because $10$ people is $2.5$ times as many as $4$ people, so the amount should be about $2.5$ times $1.5$ cups. The printer problem matches ``About $30$ minutes'' because $450$ pages at about $15$ pages per minute takes about $30$ minutes. The paint comparison matches ``Large can is a little cheaper'' because the estimated cost per gallon is lower.}

\worksheetHeader{Reasonableness Check Patrol}

\begin{errorDetective}{Does the Answer Make Sense?}

\bigskip

Problem: $210$ miles on $6$ gallons. How far on $10$ gallons?
\errorWork{Estimate about $350$ miles. Exact answer: $35$ miles. ``That looks fine.''}
Correct or wrong? \answerBlank[2cm] \quad Why? \answerBlank[6cm]

\bigskip

Problem: $5$ shirts cost $\$42.50$. How much for $8$ shirts?
\errorWork{Estimate about $\$68$. Exact answer: $\$680$. ``It is bigger, so it works.''}
Correct or wrong? \answerBlank[2cm] \quad Why? \answerBlank[6cm]

\bigskip

Problem: Small paint costs about $\$24$ per gallon; large paint costs about $\$22$ per gallon.
\errorWork{The small paint is the better deal because $24$ is greater than $22$.}
Correct or wrong? \answerBlank[2cm] \quad Why? \answerBlank[6cm]
\end{errorDetective}
\worksheetAnswer{All three are wrong. (1)~An exact answer of $35$ miles is not reasonable because the estimate was about $350$ miles, and using more gas should mean traveling farther, not less. The correct exact answer is $350$ miles. (2)~$\$680$ is not reasonable because it is ten times too large compared with the estimate of about $\$68$. A misplaced decimal often causes this kind of mistake. The correct answer is $\$68$. (3)~For prices, the smaller dollar amount is the better deal. Since $22 < 24$, the large paint can costs less per gallon and is the better buy.}

\worksheetHeader{Estimate the Fundraiser}

\begin{mathScenario}{Use an Estimate to Check the Total}

The class is selling snack boxes for a fundraiser.

\bigskip
\scenarioItem{12 boxes}{$\$31.20$ total}
\scenarioItem{Goal order}{$35$ boxes}

\bigskip

Make a quick estimate for the total cost of $35$ boxes. \answerBlank[5cm]

\bigskip

Find the exact total cost. \answerBlank[5cm]

\bigskip

Explain why your exact answer is reasonable. \answerBlank[8cm]
\end{mathScenario}
\worksheetAnswer{(1)~A quick estimate is about $\$3$ per box because $31.20 \div 12 = 2.60$, which is close to $3$. Then $35 \times 3 \approx \$105$. A closer estimate is $35 \times 2.5 = 87.5$, so the exact answer should be somewhere near $\$90$ to $\$105$. (2)~The exact unit price is $31.20 \div 12 = 2.60$ dollars per box. Then $35 \times 2.60 = \$91.00$. (3)~$\$91$ is reasonable because it falls inside the estimate range and is close to the closer estimate using $\$2.50$ to $\$3$ per box.}

\worksheetHeader{Explain Why Estimation Matters}

\begin{explainIt}{Why Not Just Solve Exactly?}

Write a paragraph explaining why estimating before solving can help you avoid mistakes in proportional reasoning.

\writeLines{6}

\bigskip

Describe one example from this chapter where an estimate would quickly show that an exact answer is unreasonable.

\writeLines{6}
\end{explainIt}
\worksheetAnswer{(1)~Estimating before solving helps because it gives a target range for the answer. If the exact result ends up far outside that range, it signals a mistake such as using the wrong operation, flipping a ratio, or misplacing a decimal. Estimation also helps students think about whether the answer should be bigger or smaller based on the context. (2)~For example, if $5$ shirts cost $\$42.50$, then each shirt costs a little more than $\$8$, so $8$ shirts should cost around $\$64$ to $\$72$. If someone gets $\$680$ or $\$6.80$, the estimate immediately shows the answer is unreasonable.}

\worksheetDone